%%%%%%%%%%%%%%%%%%%%%%%%%%%%%%%%%%%%%%%%%%%%%%%%%%%%%%%%%%%%%
%% Examtask 1: Step-up converter with output filter %%
%%%%%%%%%%%%%%%%%%%%%%%%%%%%%%%%%%%%%%%%%%%%%%%%%%%%%%%%%%%%%

\task{Step-down converter}

%%%%%%%%%%%%%%%%%%%%%%%%%%%%%%%%%%%%%%%%%%%%%%
\taskGerman{Tiefsetzsteller}

A drive system is powered by $\SI{60}{\volt}$ DC.
The speed, position, and current sensors in the drive system require stable supply voltages
of $\SI{5}{\volt}$. To transform the voltage down as efficiently as possible, a buck converter is used.

\begin{germanblock}
Ein Antriebssystem wird mit $\SI{60}{\volt}$ Gleichstrom versorgt.
Die Drehzahl-, Positions- und Stromsensoren im Antriebssystem benötigen stabile Versorgungsspannungen
von $\SI{5}{\volt}$. Um die Spannung möglichst effizient herunterzutranformieren, wird ein Tiefsetzsteller verwendet.
\end{germanblock}

% figure figStepDownConverterOutputFilter 
\input{fig/Task01/figStepDownConverterOutputFilter}

\begin{table}[ht]
    \centering  % Zentriert die Tabelle
    \begin{tabular}{llll}
        \toprule
        \multicolumn{2}{l}{\textbf{General parameters:}} & \multicolumn{2}{l}{\textbf{IGBT:}} \\ 
        Input voltage: &  $U_{\mathrm{1}} = \SI{60}{\volt}$ & Collector-emitter voltage: & $u_{\mathrm{on},\mathrm{CE}} = \SI{2}{\volt}$ \\
        Output voltage: & $U_2 = \SI{5.1}{\volt}$  & &  \\
        Output power: &$\SI{2.5}{\watt} \le p \le \SI{15}{\watt}$  & \multicolumn{2}{l}{\textbf{Diode:}} \\
        Switching frequency: & $f_\mathrm{s} = \SI{100}{\kilo\hertz}$  & Forward voltage: &  $u_{\mathrm{fw}} = \SI{0.8}{\volt}$ \\
        \bottomrule
    \end{tabular}
    \caption{Parameters of the circuit.}  % Beschriftung der Tabelle
    \label{table:ex01_Parameters of the circuit}
\end{table}


\subtask{What duty cycle must be set, assuming ideal components?}{0}
\subtaskGerman{Welches Tastverhältnis muss eingestellt werden unter der Annahme idealer Komponenten?}

\subtask{What duty cycle is required when the voltage drops across the transistor and the diode are taken into account?}{0}
\subtaskGerman{Welches Tastverhältnis ist erforderlich, wenn die Spannungsabfälle über den Transistor und die Diode berücksichtigt werden?}

\subtask{The sensors have different power consumption depending on the operating point, so the buck converter delivers 
         a maximum power of $\SI{15}{\watt}$  and a minimum power of $\SI{2.5}{\watt}$.
        What inductance must be designed so that it does not operate in DCM mode within the specified power range, 
        reaching BCM mode at 2.5 watts?}{1}
\subtaskGerman{Die Sensoren haben je nach Arbeitspunkt eine unterschiedliche Leistungsaufnahme, so dass der Tiefsetzsteller 
für einen maximalen Leistung von $\SI{15}{\watt}$ und eine minimale Leistung von $\SI{2.5}{\watt}$ liefert.
Wie groß muss die Induktivität ausgelegt werden, damit er im angegebenen Leistungsbereich nicht im DCM-Mode arbeitet, wobei er bei $\SI{2.5}{\watt}$ den BCM-Mode erreicht.
}

The sensors are susceptible to voltage fluctuations, so the power supply must have a maximum ripple of
$\pm \SI{2.5}{\percent}$ of the output voltage. To achieve this, a filter capacitor is added.

\begin{germanblock}
Die Sensoren sind gegenüber Spannungsschwankungen anfällig, so dass die Spannungsversorgung eine maximale Welligkeit von
 $\pm \SI{2.5}{\percent}$ der Ausgangsspannung besitzen darf. Um dieses  zu erreichen, wird ein Filterkondensator eingebaut.
\end{germanblock}

\subtask{What value should be chosen for the capacitance of the filter capacitor?}{0}
\subtaskGerman{Welche Größe muss für die Kapazität des Filterkondensators gewählt werden?}

\begin{germanblock}Wie groß muss die Kapazität des Filterkondensators gewählt werden?
In einem Störfall kommt es zum Ausfall der Drehzahl- und Positionssensoren. Dadurch fällt die abgenommene Leistung auf einen Wert von $\SI{1}{\watt}$.
\end{germanblock}


\subtask{How large must the filter capacitor be so that, in the worst case, the voltage fluctuation is $\pm \SI{10}{\percent}$ ?}{0}
\subtaskGerman{ Wie hoch ist die Ausgangsspannung für diesen Störfall?}